\subsection{Fase 3: Projeto da Avaliação}
\label{sec:fase3-plano}

A avaliação será externa, em caixa-preta, com roteiro fixo e o mesmo veículo nos dois canais. O plano operacional e os instrumentos estão em \texttt{docs/fase3-*.md} e \texttt{dados/} no repositório. A Tabela~\ref{tab:fase3-procedimentos} resume a rastreabilidade entre as métricas da Fase 2 e os procedimentos. Cada sessão terá executor, registrador e conferência posterior por avaliador independente, com vínculo entre registro, contexto e evidência anonimizada.

\begin{table}[htbp]
\centering
\caption{Procedimentos e registros planejados para a avaliação.}
\label{tab:fase3-procedimentos}
\small
\begin{tabular}{p{0.08\textwidth}p{0.28\textwidth}p{0.51\textwidth}}
\toprule
Roteiro & Métricas & Instrumento e registro \\
\midrule
P1 & M1.1.1, M1.1.2 & Capturas; presença dos campos por infração. \\
P2 & M1.2.1, M1.2.2 & Capturas pareadas; classificação da consistência. \\
P3 & M1.3.1, M1.3.2 & Vídeo; passos e transições em três execuções. \\
P4 & M2.1.1, M2.1.2 & Vídeo; quadros ou timestamps e tempos. \\
P5 & M2.2.1 & DevTools e contador por UID no Android; bytes. \\
P6 & M3.1.1, M3.1.2 & Três tentativas por contexto; registro de conclusão. \\
P7 & M3.2.1--M3.2.3 & Vídeo e registro de tentativas por canal e rede. \\
\bottomrule
\end{tabular}
\par\smallskip{\footnotesize Fonte: elaborado para o plano de avaliação.}
\end{table}

O contexto de referência será o Portal de Serviços no Chrome em computador e o aplicativo em Android intermediário, ambos em Wi-Fi. Para cobertura, acrescentam-se Firefox no computador, Chrome no Android, Android de entrada e iOS se houver empréstimo. Modelos, memória, sistema, versões instaladas e taxa de gravação serão registrados por sessão; esses recursos ainda não foram confirmados. Sem iPhone, o iOS será não avaliado e as conclusões do aplicativo se restringirão ao Android.

As consultas começam após autenticação pelo titular. Mantém-se o cache nos dois canais e descarta-se um aquecimento por bloco, sem presumir equivalência interna dos caches. P4 e P5 terão dez medidas válidas por canal, com limite de vinte tentativas; P7 terá dez tentativas por canal e condição, preservando falhas. O limite de espera será 60~s. Os tempos serão resumidos por média, mediana e P90 pelo posto mais próximo; o julgamento usa a mediana. P6 considera concluído o contexto com sucesso em pelo menos duas de três tentativas.

As condições de rede serão Wi-Fi, saída 4G e perfil degradado aplicado no ponto de acesso Linux aos dois clientes: atraso de 300~ms com variação de 50~ms, banda de 1~Mbit/s e perda de 1\% por sentido. O piloto deve validar interfaces, configuração e instrumentos. O volume só permitirá comparação direta se os contadores medirem escopos equivalentes; caso contrário, será reportado separadamente, com limitação explícita. Métricas indisponíveis não recebem valor zero nem julgamento favorável.

A amostra intencional prevê até três veículos, com mínimo de um veículo com infração e consentimento escrito. Busca-se diversidade de situações, órgãos e infrações de velocidade, sem pressupor sua disponibilidade. Selecionam-se previamente até três infrações por veículo para P1/P2; P3--P7 mantêm um veículo de referência e a mesma infração de detalhamento, sem tratar repetições como casos independentes. A mesma infração será comparada em intervalo de até dez minutos; pares fora dessa janela serão preservados como inelegíveis e refeitos. Termos, correspondências, vídeos originais e HAR permanecerão privados; somente códigos e evidências anonimizadas serão publicados. A gravação começa após o login. A ausência de veículo ou consentimento bloqueia a coleta.

O denominador de P2 contém os 15 campos de referência e dois campos adicionais apenas nas infrações de velocidade; ausências não são paridade. M1.2.1 e M1.2.2 descrevem consistência do par, sem indicar qual canal está correto. Aplicam-se os limites contínuos da Fase 2 aos valores sem arredondamento. O registro de sessão preservará a versão do plano, a liberação do piloto, as duas leituras da conferência e a decisão final; os arquivos CSV conterão a medida conferida e sua evidência.

Propõem-se preparação em 14/11/2026 (Hugo), piloto em 15/11 (Hugo e Philipe), coleta de referência em 16/11 (Hugo e Philipe), contextos e redes em 17/11 (Pedro Druck, Hugo e Philipe) e contingência/revisão em 18/11 (equipe de coleta; Daniel revisa). Pedro Henrique revisará o plano até 15/11. Esses marcos preservam os prazos de coleta em 18/11 e análise em 20/11; horários com titulares ainda dependem de agendamento. O plano não constitui evidência de execução: consentimentos, inventário efetivo, capturas autenticadas e aprovação do revisor permanecem pendentes.

A análise será conduzida por Pedro Henrique e Philipe em 19--20/11, com revisão de Heitor,
antes da Entrega 2 em 22/11. Os prazos constam das issues \#47 e \#48; a conferência dos campos
de data no Project permanece pendente por restrição de acesso. O quadro de bloqueios e o
cronograma operacional em \texttt{docs/fase3-cronograma.md} registram responsáveis, datas e
condições para replanejar sem omitir métricas não medidas.
